\section{Scalar-optimal actual radical partitions without reusable flags}
\label{dp-section}

This section gives an arithmetic subclass of the unchanged FCRT problem.
It neither bounds the scalar ABC defect from above nor constructs an ABC
counterexample. All divisors below are divisors of actual integer radicals.

For the primitive endpoint \((1,c-1,c)\), with
\(c=2^e3^f\) and integers \(e,f\ge2\), put
\[
 M=2^e,\quad N=3^f,\quad A=M/2,\quad B=N/3,
 \quad Q=\operatorname{rad}(c-1).
\]
The previously established exact optimization on the stratum
\[
 \max(M,N)<2\min(M,N)-1,\qquad AB>Q
\]
is
\begin{gather}
 B_{\rm FCRT}=B_{\rm SCRT}=B_{\rm PBT}
 =D+\min_{H\mid Q}\operatorname{dist}
       \bigl(\log H,[\log(Q/B),\log A]\bigr),\label{dp-exact}\\
 D=\log(AB/Q)=\log\!\left(\frac{c}{\operatorname{rad}(c(c-1))}\right)>0.
 \label{dp-defect}
\end{gather}
Here no nonempty proper compatible face exists at either source, and
every sink prime is assigned with its entire power in the endpoint.
We recall the two inherited facts needed for this reduction.

If \(p^e\Vert c\), \(m=p^e\), and \(c-1=du\) with \(d,u>1\)
and \(u\equiv-1\pmod m\), then \(c\ge m(2m-1)\).
Indeed, write \(c=mC\), \(d=mt+1\), \(u=ms-1\) with positive
integers \(s,t\). Exactness of \(p^e\Vert c\) excludes \(C=m\).
Expansion gives \(C+t=s(mt+1)\). If \(s=t=1\) then \(C=m\),
excluded; otherwise \(C\ge2m-1\). The balance condition above
therefore excludes every proper compatible packet for both exact source
moduli. Every compatible block must use the full sink set; a block with
both sources is unsaturated, and a singleton block cannot improve on
giving all sinks to that source. Thus the full optimization is
\[
 \log\min_{H\mid Q}\max(A/H,1)\max(BH/Q,1).
\]
Splitting into the three positions of \(\log H\) relative to the stated
interval proves \eqref{dp-exact}. These are the ordinary no-face and
optimization results used here; no uncompiled draft is reclassified as
a Lean proof.

\begin{theorem}[A complete two-factor partition]\label{dp-partition}
Let \(u,v\ge1\) and \(k\ge2\) be integers, set \(x=2^u3^v\),
\(e=2u\), \(f=2v\), and suppose
\begin{equation}\label{dp-hypotheses}
 x>6,\qquad \frac34\le\rho:=\frac{2^u}{3^v}\le\frac43,
 \qquad 7^k\mid x-1,\quad 5^k\mid x+1.
\end{equation}
Then the hypotheses of \eqref{dp-exact} hold, all proper flags are absent,
and the actual coprime radicals
\[
 H=\operatorname{rad}(x-1),\qquad
 J=\operatorname{rad}(x+1),\qquad Q=HJ
\]
satisfy \(H\le A\) and \(J\le B\). Consequently
\begin{equation}\label{dp-scalar-optimum}
 B_{\rm FCRT}=D>(k-1)\log35-\log6.
\end{equation}
\end{theorem}
\begin{proof}
Since \(x\) is even, \(x-1,x+1\) are odd and coprime. Their radicals
therefore multiply to exactly \(Q=\operatorname{rad}(x^2-1)\).
For any positive integer \(m\) with \(p^k\mid m\),
\(\operatorname{rad}(m)\le m/p^{k-1}\), regardless of all other
prime powers. Hence
\[
 H\le\frac{x-1}{7^{k-1}}<\frac{x}{7},\qquad
 J\le\frac{x+1}{5^{k-1}}\le\frac{x+1}{5}\le\frac{x}{4}.
\]
The last inequality uses \(x\ge4\). Meanwhile
\[
 A=\frac{x\rho}{2}\ge\frac{3x}{8},\qquad
 B=\frac{x}{3\rho}\ge\frac{x}{4}.
\]
Thus \(Q/B\le H\le A\). The first packet uses all prime powers of
\(x-1\), and the second all prime powers of \(x+1\). They are disjoint
and exhaust every sink, with no duplicated credit or putative subface.

For the no-flag hypothesis, \(M=x\rho\), \(N=x/\rho\),
\(\max(M,N)/\min(M,N)\le16/9\), and \(\min(M,N)\ge3x/4\).
Therefore
\[
 2\min(M,N)-\max(M,N)
 \ge\frac29\min(M,N)\ge\frac{x}{6}>1.
\]
The powers \(M,N\) are exact full prime powers in \(c\), so the
no-face result applies. Moreover
\[
 Q\le\frac{x^2-1}{35^{k-1}}<\frac{x^2}{6}=AB.
\]
The exhibited divisor makes the distance in \eqref{dp-exact} zero, while
\[
 \frac{AB}{Q}\ge
 \frac{35^{k-1}x^2}{6(x^2-1)}>\frac{35^{k-1}}6.
\]
Taking logarithms proves \eqref{dp-scalar-optimum}. No factorization
of the remaining parts of \(x-1\) or \(x+1\) was required.
\end{proof}

\begin{lemma}[Simultaneous actual residues]\label{dp-residues}
For every integer \(k\ge2\), put \(L=12\cdot35^{k-1}\).
For integers \(r\ge0\), \(s\ge1\), define
\begin{equation}\label{dp-exponents}
 u=Lr+L/2=6\cdot35^{k-1}(2r+1),\qquad v=Ls.
\end{equation}
Then \(x=2^u3^v\) satisfies \(7^k\mid x-1\) and \(5^k\mid x+1\).
\end{lemma}
\begin{proof}
The binomial induction
\(z\equiv1\pmod p\Rightarrow z^{p^{k-1}}\equiv1\pmod{p^k}\)
gives
\[
 2^{3\cdot7^{k-1}}\equiv1\pmod{7^k},\qquad
 2^{2\cdot5^{k-1}}\equiv-1\pmod{5^k}.
\]
For the second congruence apply it to \(-4\equiv1\pmod5\) and use
that \(5^{k-1}\) is odd. The exponent \(u\) is a multiple of
\(3\cdot7^{k-1}\), while
\[
 \frac{u}{2\cdot5^{k-1}}=3\cdot7^{k-1}(2r+1)
\]
is odd. Thus \(2^u\) has the required signs at both moduli.
The same induction applied to \(3^6\equiv1\pmod7\) and
\(3^4\equiv1\pmod5\), together with the respective divisibilities
\(6\cdot7^{k-1}\mid v\) and \(4\cdot5^{k-1}\mid v\), gives
\(3^v\equiv1\) at both moduli. Multiplication finishes the proof.
\end{proof}

\begin{theorem}[Infinitely many scalar-optimal no-flag endpoints]
\label{dp-infinite}
For every fixed integer \(k\ge2\), infinitely many pairs \(r,s\) in
\eqref{dp-exponents} satisfy \eqref{dp-hypotheses}; they may be chosen
with \(\rho\to1\) and \(x\to\infty\). Across these fixed-\(k\)
families, the absolute scalar defect \(D\) is unbounded, while the
radical partition penalty in \eqref{dp-exact} is identically zero.
\end{theorem}
\begin{proof}
Let \(\alpha=\log3/\log2\), irrational by unique factorization.
Pigeonhole approximation gives positive integers \(q_j\) and integers
\(a_j\) with \(\beta_j=|q_j\alpha-a_j|>0\) tending to zero.
Choose a positive integer \(h_j\) nearest to \(1/(2\beta_j)\).
Then \(h_j\beta_j\to1/2\). For \(s_j=h_jq_j\), the fractional
part of \(\alpha s_j\) tends to \(1/2\), whether the approximation
error has positive or negative sign. Also \(s_j\to\infty\).
Take \(r_j=\lfloor\alpha s_j\rfloor\), eventually positive. For
the fixed number \(L=12\cdot35^{k-1}\),
\[
 r_j+\tfrac12-\alpha s_j\longrightarrow0,
 \qquad
 \log\rho_j=L\log2\,(r_j+\tfrac12-\alpha s_j)\longrightarrow0.
\]
Thus the balance conditions and \(x>6\) hold eventually, while
Lemma~\ref{dp-residues} supplies both divisibilities.
Theorem~\ref{dp-partition} applies to all those actual inputs.

For any real threshold \(T\), first choose \(k\) so that
\((k-1)\log35-\log6>T\), and then take arbitrarily large members of
that fixed-\(k\) family. Their defects exceed \(T\). A diagonal choice
makes both \(c\) and \(D\) increase without bound. The radical
partition penalty stays zero throughout.
\end{proof}

This disproves the particular added claim that absence of proper flags
and sufficiently large absolute defect force a positive partition penalty.
It does not disprove FCRT-1 or SCRT-0. In particular, no fixed
\(\varepsilon>0\) is exhibited for which the ABC excess
\(\log c-(1+\varepsilon)\log\operatorname{rad}(c(c-1))\)
is unbounded. Nor is an upper bound for that expression, for \(D\),
or for the partition penalty on all other endpoints proved here.

The global maximum consecutive divisor gap of \(Q\) can be much less
relevant than the required central interval. The present proof supplies
one actual divisor in that interval, without estimating the maximum gap
over the entire divisor set. All unspecified prime factors remain in
their full integer packets. The surviving general problem includes both
the scalar defect and actual central divisor gaps when such witnesses
are unavailable.

\paragraph{Ordinary and finite verification scope.}
The complete ordinary argument received independent internal review.
The standard-library program
\texttt{replay\_\allowbreak divisor\_\allowbreak partition.py --check}
checks three actual balanced inputs by exact integer inequalities and
the prescribed prime-power divisibilities, together with 66 residue-only
instances for \(k=2,\ldots,12\). The latter are not asserted to be balanced.
Its canonical certificate SHA-256 is
\texttt{6b7a363b91b7fad5\allowbreak c570696ed8b5fcd5\allowbreak
0be50d73607a7d0c\allowbreak 0d829bc1b75c551d}.
It neither factors the remaining neighbours nor computes their radicals
or logarithms, and it does not prove the infinite-family theorem.
No new Lean theorem is claimed in this section.
